%% ===================================================================
\section{Running coupling and DGLAP}\label{sec:rc}
%% ===================================================================

Throughout this section $d_\perp=2-\epsilon$, so that in the standard notation $D=4-2\epsilon_{\rm std}$ one has $\frac2\epsilon=\frac{1}{\epsilon_{\rm std}}$. We set $n_f=0$, so $\beta_0=\frac{11}{3}N_c$, and $\alpha_s=g^2/4\pi$.

\subsection{The splitting function}\label{sec:Phat}

The real part of the gluon splitting function~\cite{Gribov:1972ri,Altarelli:1977zs,Dokshitzer:1977sg} appears in every term that contains the splitting of one gluon into two. In the form in which it arises here,
\begin{equation}
\Phat(\xi)\equiv\xi(1-\xi)\frac{d_\perp}{2}+\frac{\xi}{1-\xi}+\frac{1-\xi}{\xi},\qquad \Phat(\xi)\Big|_{d_\perp=2}=\frac{P^{\rm real}_{gg}(\xi)}{2N_c},
\label{eq:Phat}
\end{equation}
where $P_{gg}(\xi)=2N_c\big[\frac{\xi}{(1-\xi)_+}+\frac{1-\xi}{\xi}+\xi(1-\xi)\big]+\frac{\beta_0}{2}\delta(1-\xi)$. The key identity is
\begin{equation}
\int_{\Lambda/k^+}^{1-\Lambda/k^+}d\xi\,\Phat(\xi)\Big|_{d_\perp=2}=2\log\frac{k^+}{\Lambda}-\frac{11}{6}+\Ord\Big(\frac{\Lambda}{k^+}\Big),\qquad\text{i.e.}\qquad \frac{\beta_0}{N_c}=\frac{11}{3}=-2\Big[\int d\xi\,\Phat-2\,\delta Y_T\Big].
\label{eq:intPhat}
\end{equation}
The same function therefore produces a rapidity logarithm (from its endpoints) and the $\beta_0$ coefficient (from its integral).

\subsection{\texorpdfstring{$\beta_0$}{beta0} from the gluon self-energy before the target (G1-II)}\label{sec:b0II}

In the self-energy of the emitted gluon, Eq.~\eqref{eq:grho}, the $\xi$-integrand of the gluon loop is $\big[-2+\xi(1-\xi)+\frac{1}{\xi(1-\xi)}\big]$. This is identically $\Phat(\xi)$ at $d_\perp=2$, because $\frac{1-\xi}{\xi}+\frac{\xi}{1-\xi}=\frac{1}{\xi(1-\xi)}-2$. By Eq.~\eqref{eq:intPhat} the coefficient $\frac{11}{3}+4\log\frac{\Lambda}{k^+}$ of the UV pole of the self-energy is exactly $-2\int d\xi\,\Phat$. The other one-loop diagrams of Eq.~\eqref{eq:grho}, in which the loop is attached to the valence charge, add $2\delta Y_T+\lV$ to this coefficient, which gives the total $\frac{11}{3}-2\delta Y_T+\lV$ of Eq.~\eqref{eq:grhologs}. Relative to LO, Eq.~\eqref{eq:finII1} together with its subtracted part is, per ordering,
\begin{equation}
\frac{\text{G1-II(1)}}{\rm LO}=-\frac{\alpha_sN_c}{4\pi}\Big\{\Big[-\frac2\epsilon-\log\frac{r^2\bar\mu^2e^{2\gamma_E}}{4}\Big]\Big(\frac{11}{3}-2\delta Y_T+\lV\Big)-\frac{67}{9}+\frac{\pi^2}{3}+\dots\Big\},\qquad r=|x-w|,
\end{equation}
and its $\beta_0$ part is
\begin{equation}
\frac{\alpha_s\beta_0}{4\pi}\Big[\frac{1}{\epsilon_{\rm std}}+\log\frac{r^2\bar\mu^2e^{2\gamma_E}}{4}\Big]\times{\rm LO}\qquad\text{per ordering (with $r'=|x'-w'|$ in the conjugate).}
\label{eq:b0II}
\end{equation}

\subsection{\texorpdfstring{$\beta_0$}{beta0} from the loop across the target (G1-V, part 2b)}\label{sec:b0V}

In G1-V the two gluons produced before the target recombine after it. The colour operator $\Ord_{\rm V2b}$ does not vanish when the two gluons coincide, and the row is UV divergent. To extract the coefficient, set $x=z+\mathbf r$ and let $\mathbf r\to0$. Then $\mathcal D\to k^+(y-z)^2$, the phase $e^{-ik\cdot(z-x)p^+/k^+}\to1$, and $\Ord_{\rm V2b}\to-N_c\Ord_{\rm LO}$ for a gluon at $z$ with source $y$. Averaged over the direction of $\mathbf r$, the six brackets of Eq.~\eqref{eq:B16} give, with $X=x'-w'$, $Y=y-z$, $\xi=p^+/k^+$,
\begin{align}
&\frac{X^i(z-x)^m}{(z-x)^2}\Big[\xi(1-\xi)\mathcal B_1-\tfrac{\xi}{1-\xi}\mathcal B_2+(1-\xi)\mathcal B_3-\tfrac{1-\xi}{\xi}\mathcal B_4+\xi\mathcal B_5-\mathcal B_6\Big]^{im}\notag\\
&\qquad\longrightarrow\frac{X\cdot Y}{r^2}\Big[\frac{d_\perp}{2}\xi(1-\xi)+\frac{\xi}{1-\xi}-(1-\xi)+\frac{1-\xi}{\xi}-\xi+1\Big]=\frac{X\cdot Y}{r^2}\,\Phat(\xi).
\end{align}
The six terms, in the order of the brackets, are $\frac{d_\perp}{2}\xi(1-\xi)$, $\frac{\xi}{1-\xi}$, $-(1-\xi)$, $\frac{1-\xi}{\xi}$, $-\xi$ and $+1$. With $\int dp^+=k^+\int d\xi$,
\begin{equation}
\frac{d^3N_{\rm V2b}}{d^3k}\bigg|_{\rm UV}=\frac{\alpha_sN_c}{2\pi}\Big[\int d\xi\,\Phat(\xi)\Big]\log\frac{R^2}{r_0^2}\times{\rm LO}=\frac{\alpha_sN_c}{2\pi}\Big[2\delta Y_T-\frac{11}{6}\Big]\log\frac{R^2}{r_0^2}\times{\rm LO},
\label{eq:2bUV}
\end{equation}
where LO refers to a gluon at $z$ with source $y$, $r_0$ is a short-distance cutoff ($\log\frac{1}{r_0^2}\leftrightarrow\frac{1}{\epsilon_{\rm std}}+\log\bar\mu^2$), and the upper cutoff $R$ of $|x-z|$ is set by the kernel ($|x-z|\lesssim|y-z|$) and by the phase ($|x-z|\lesssim1/(\xi|\mathbf k|)$).
\begin{itemize}
\item The $2\delta Y_T$ part is the UV-divergent form of a JIMWLK virtual term. It sits in the endpoint terms (the subtracted part of G1-V(2b)).
\item The $-\frac{11}{6}$ part is a $\beta_0$ term, $-\frac{\alpha_s\beta_0}{4\pi}\log\frac{R^2}{r_0^2}\times{\rm LO}$.
\end{itemize}
In the same cutoff language Eq.~\eqref{eq:b0II} reads $+\frac{\alpha_s\beta_0}{4\pi}\log\frac{r^2}{r_0^2}\times$LO. The loop before the target and the loop across the target therefore come with \emph{opposite} signs.

\subsection{Renormalization-group bookkeeping}\label{sec:RG}

With the bare coupling, ${\rm LO}(\alpha_0)={\rm LO}(\alpha_s(\mu))\big[1-\frac{\alpha_s\beta_0}{4\pi}\frac{1}{\epsilon_{\rm std}}\big]$. For the NLO cross section to be finite and independent of $\mu$, the UV poles of all $\beta_0$ terms must add up to $+\frac{\alpha_s\beta_0}{4\pi}\frac{1}{\epsilon_{\rm std}}\times$LO.
\begin{itemize}
\item G1-II with its complex conjugate gives $2\times\frac{\alpha_s\beta_0}{4\pi}\frac{1}{\epsilon_{\rm std}}$, twice the required amount.
\item G1-V(2b) gives $-\frac{\alpha_s\beta_0}{4\pi}\frac{1}{\epsilon_{\rm std}}$ per ordering.
\end{itemize}
The balance therefore holds if G1-V(2b) enters with net weight one relative to one ordering of G1-II.
\note{This has to be checked against the factor 2 in Eq.~\eqref{eq:VV} and the complex conjugate; it is the main open normalization question for the running coupling.}
%TODO: verify the net weight of G1-V(2b) relative to G1-II (beta0 pole balance).
If it holds, the surviving $\beta_0$ logarithms are
\begin{equation}
\frac{\alpha_s\beta_0}{4\pi}\Big[\log\big(r^2\mu^2\big)+\log\big(r'^2\mu^2\big)-\log\big(R^2\mu^2\big)\Big]\times{\rm LO}+\text{const}\times{\rm LO}\quad\Longleftrightarrow\quad\alpha_s(\mu)\ \to\ \frac{\alpha_s(1/r^2)\,\alpha_s(1/r'^2)}{\alpha_s(1/R^2)}.
\label{eq:rcstruct}
\end{equation}
Its $\mu$ dependence cancels that of $\alpha_s(\mu)$ in LO. Eq.~\eqref{eq:rcstruct} is a ratio of three running couplings, of the same type as the ``triumvirate'' structures found in running-coupling corrections to small-$x$ evolution~\cite{Balitsky:2006wa,Kovchegov:2006vj} and to the LO gluon spectrum~\cite{Horowitz:2010yg}. The two couplings in the numerator sit at the emission scales of the projectile in the amplitude ($r=|x-w|$) and in its conjugate ($r'=|x'-w'|$). The coupling in the denominator sits at the final-state scale, $R\sim1/|\mathbf k|$ for $|\mathbf k|\gg1/|y-z|$. Couplings associated with the target are contained in the Wilson-line correlators here.

\subsection{Not all \texorpdfstring{$\beta_0$}{beta0} terms are running coupling}\label{sec:KLSZ}

The logarithm $\log(r^2\mu^2)$ of Eq.~\eqref{eq:rcstruct} is evaluated at the projectile emission size $r=|x-w|$. In the LO spectrum $r$ ranges up to the colour correlation length $R_p$ of the projectile, i.e.\ the range of $\langle\rho\rho\rangle$. The target cuts the emission off at $1/Q_s$ through the factor $U(w)-U(x)$. Following Kovner, Lublinsky, Skokov and Zhao~\cite{Kovner:2023vsy} (see also Ref.~\cite{Altinoluk:2023hfz}):
\begin{itemize}
\item when $R_p\gg1/Q_s$, the $\beta_0$ logarithm is $\log(R_p^2Q_s^2)$. It is large, it is \emph{not} a running-coupling logarithm, and it must be resummed by DGLAP evolution of the projectile from $1/R_p$ to $Q_s$;
\item only when $R_p\lesssim1/Q_s$ is it a running-coupling logarithm.
\end{itemize}
In our calculation this distinction appears as follows. The $\beta_0$ term of G1-II is the $\delta(1-\xi)$ (virtual) part of $P_{gg}$. The real parts $\frac{\xi}{1-\xi}+\frac{1-\xi}{\xi}+\xi(1-\xi)$, with the same collinear logarithm, are in the initial-state rows discussed next. Together they form $P_{gg}\otimes{\rm LO}$, as required by DGLAP evolution of the projectile. The $\beta_0$ terms and the DGLAP terms therefore cannot be separated row by row.

\subsection{The constant \texorpdfstring{$-\frac{67}{9}+\frac{\pi^2}{3}$}{-67/9+pi2/3}}\label{sec:CMW}

$-\frac{67}{9}+\frac{\pi^2}{3}=-2\big(\frac{67}{18}-\frac{\pi^2}{6}\big)=-\frac{2K}{C_A}$, where $K=C_A\big(\frac{67}{18}-\frac{\pi^2}{6}\big)-\frac59n_f$ is the Catani--Marchesini--Webber constant~\cite{Catani:1990rr}. Per ordering, G1-II gives $+\frac{\alpha_sK}{2\pi}\times$LO. In the CMW (bremsstrahlung) scheme, $\alpha_s\to\alpha_s\big(1+\frac{K\alpha_s}{2\pi}\big)$ absorbs exactly this term per emission. With the complex conjugate G1-II gives twice this, the same factor 2 as for the UV pole. Whether the remainder is supplied by the finite constant of G1-V(2b) can be decided only once part 2b is completed in dimensional regularization.

\subsection{DGLAP: collinear limits of the real terms}\label{sec:dglap}

The splitting function~\eqref{eq:Phat} appears in all terms with a gluon splitting, see Table~\ref{tab:dglap}. All these limits were checked numerically (App.~\ref{app:checks}).

\begin{table}[!tbp]
\centering
\renewcommand{\arraystretch}{1.25}
\begin{tabular}{lL{4.4cm}L{8.2cm}}
\toprule
Row & splitting & limit ($\xi$: momentum fraction of the measured gluon in the splitting)\\
\midrule
G1-II & virtual, before the target & loop integrand $=\Phat(\xi)$ exactly; gives $\beta_0$ and $\delta Y_T$\\
G1-V(2b) & virtual, across the target & UV coefficient $=\Phat(\xi)$, $\xi=p^+/k^+$; Eq.~\eqref{eq:2bUV}\\
G2-I(2) & real, initial state & for $|z-w|\ll|x-z|$: integrand $\to\frac{(x-z)\cdot(x'-z)}{(x-z)^2(x'-z)^2}\frac{\Phat(\xi)}{|z-w|^2}$, $\xi=\frac{k^+}{k^++p^+}$\\
G3-III(2) & interference of initial and final state & same limit with coefficient $-1$ per ordering ($-2$ with c.c.)\\
G2-IV & real, final state (fragmentation) & same limit with coefficient $+1$; $\mathcal P_{gg}$ and $D_{g/g}$\\
\bottomrule
\end{tabular}
\caption{The gluon splitting function in the NLO cross section.}
\label{tab:dglap}
\end{table}

\paragraph*{Collinear limit of the real rows.} Put the measured gluon and the unobserved gluon close together, $|z-w|=|z-w'|=r\ll|x-z|,|x'-z|$, and average over the direction of $z-w$. With the common prefactor $\frac{1}{(2\pi)^3}\frac{g^4}{4\pi^4}\int\frac{dp^+}{2\pi}$, all three real rows reduce to
\begin{equation}
c_{\rm row}\,\frac{\Phat(\xi)}{(k^++p^+)^2}\,\frac{(x-z)\cdot(x'-z)}{(x-z)^2(x'-z)^2}\,\frac{1}{r^2}\times(\text{colour operator at }w=w'=z),
\label{eq:collinear}
\end{equation}
with $c_{\rm II.I(2)}=+1$, $c_{\rm III.III(2)}=-1$ (per ordering) and $c_{\rm II.IV}=+1$. The kinematic factor is the LO kernel of the parent gluon at $z$, with momentum $k^++p^+$, times $\Phat(\xi)/r^2$. Since $\int dp^+/(k^++p^+)^2=\frac{1}{k^+}\int d\xi$, this is the DGLAP kernel times the collinear logarithm $\int dr^2/r^2$, convoluted with the LO spectrum of the parent.
\begin{itemize}
\item The weights $1:-2:1$ (with the complex conjugate) are those of $|\psi_{\rm initial}-\psi_{\rm final}|^2$. When the target cannot resolve the pair ($r\ll1/Q_s$, so that $U(w)\simeq U(z)$), splitting before and after the target are indistinguishable, and the collinear singularity at $r\to0$ cancels among the real rows. This requires the colour operators of the three rows at $w=w'=z$ to coincide.
\note{With a general adjoint $U$ and the index orders of the draft notes, the three colour operators at $w=w'=z$ do not coincide; the index placement in G2-I(2), G2-IV and G3-III(2) must be made consistent before this cancellation can be verified.}
%TODO: make the Wilson-line index placement in G2-I(2), G2-IV, G3-III(2) consistent; verify collinear cancellation.
\item What survives is the region in which the target resolves the pair, $1/Q_s\lesssim r\lesssim|x-z|$. There the initial-state rows give $P_{gg}\otimes{\rm LO}$ times $\log(R_p^2Q_s^2)$, the projectile DGLAP logarithm of Sec.~\ref{sec:KLSZ}, and the final-state row gives the fragmentation logarithm.
\end{itemize}

\paragraph*{Final state.} In G2-IV the $z$ integration produces a collinear pole, which is absorbed into the gluon fragmentation function, $D_{g/g}(\xi)=\delta(1-\xi)-\frac{\alpha_s}{2\pi}\frac{1}{\bar\epsilon}\mathcal P_{gg}(\xi)$. This is the standard $\MSbar$ renormalization of $D_{g/g}$. Its virtual counterpart, the $\frac{\beta_0}{2}\delta(1-\xi)$ of $P_{gg}$, comes from the self-energies of Secs.~\ref{sec:b0II} and~\ref{sec:b0V}. The remaining finite terms of G2-IV are its finite part, Sec.~\ref{sec:finiteG2}.

\paragraph*{Scheme.} The $d_\perp$ dependence of $\Phat$ in Eq.~\eqref{eq:Phat} sits only in the $\xi(1-\xi)$ term, as it arises from the trace $\delta^{ij}\delta^{ij}=d_\perp$ in our expressions. When the $1/\epsilon$ poles of G1-II and G1-V(2b) are combined, the $\Ord(\epsilon)$ part of $\int d\xi\,\Phat$ produces finite constants, and the two rows must be evaluated in the same scheme.
